\chapter{Quantum Mechanics}
\label{chap:quantum-mechanics}

% Closed question: How can information move between histories while preserving
% admissibility?

\begin{chapteressay}

A photon enters a Mach-Zehnder interferometer. A beam splitter sends the
amplitude into two arms. Mirrors guide the arms back together. A detector waits
at the end.

If no one asks which arm the photon used, the two arms interfere. The output
pattern depends on both alternatives. If a which-path detector is placed in one
arm, the interference disappears. The history has not merely become known. The
record has changed.

This is the chapter's governing example. Classical mechanics selected one
admissible continuation. Quantum mechanics can keep more than one continuation
admissible until a witnessing interaction turns the alternatives into a record.
The hard point is not that the photon is mysterious. The hard point is that
``after'' is cheap in prose and expensive in physics. A sentence can say that
the photon went through the splitter and then through one arm and then to the
detector. The apparatus refuses that grammar unless the corresponding records
exist.

The speedometer has an analogy at its resolution floor. If the wheel sensor has
one-tooth resolution, events inside one tooth interval are not available as
separate records. The instrument may carry a distribution of compatible local
states, but it may not pretend to know a finer history. Quantum mechanics makes
that humility fundamental. Below the scale carried by the measurement, path
distinctions are not merely unknown. They are not yet admissible facts.

The same structure appears in a sealed routing room. Two diplomatic pouches
enter with correlated labels. The room can move them, sort them, and deliver
them. But opening a pouch changes the chain of custody. A correlation that was
carried silently becomes a witnessed record only when the seal is broken and
the two pouches are compared. The room did not transmit a message faster than
light. It preserved a joint constraint until measurement made the constraint
visible.

The code has been preparing this language. \lean{Science} can be repeatable,
hypothetical, or theoretical. \lean{Knowledge} grows by ledger entries.
\lean{TRUTH} tests whether knowledge can be carried in both directions.
\lean{WITNESSED} introduces a witness and asks for a reciprocal gospel order.
The comedy of the names should not hide the structure: knowledge is not free
possession. It is a carried process with a record.

The closed question is:

\begin{center}
\emph{How can information move between histories while preserving
admissibility?}
\end{center}

The answer is that quantum information is transported as a constrained
record, not as a classical fact waiting underneath. Interference keeps
alternative histories jointly admissible. Decoherence occurs when an
environment carries which-history information. Entanglement preserves a joint
constraint across separated measurements. A witness closes the record.

Chapter 7 therefore makes history plural until measurement makes it singular.
The book does not need the full Hilbert-space formalism here. It needs the
measurement grammar: admissible alternatives, transport, witnessing, and the
physical cost of turning a possible distinction into a recorded one.

\end{chapteressay}

% Phenomenon arcs use: 1/3 setup -> phenomenon box -> 2/3 structural interpretation.

\section{The Wittgenstein Effect: After is Free}
\label{se:the-wittgenstein-effect}

Language lets us put ``after'' between almost anything. The photon passes the
beam splitter after it enters. It reaches the detector after traveling. A
sentence can line events up without paying for the records that distinguish
them.

The Mach-Zehnder interferometer exposes the cost. With both arms open and no
which-path detection, the two alternatives recombine. The detector statistics
come from interference between amplitudes. If a detector is installed in one
arm, the apparatus has created a new record. The interference pattern changes.

\begin{phenom}{Wittgenstein Effect}
\label{ph:wittgenstein-effect}

\PhStatement Grammar can order events more cheaply than physics can record
them; a claimed ``after'' becomes physical only when the apparatus carries the
distinction.

\PhOrigin Interferometer experiments show that path language and path records
come apart whenever alternatives remain coherent.

\PhObservation In a Mach-Zehnder interferometer, path alternatives interfere
when no which-path record exists. Adding a detector changes the output
statistics by making the path distinction physical.

\PhConstraint The measurement may not promote a syntactic sequence into a
physical sequence unless the apparatus carries the relevant record.

\PhConsequence Quantum mechanics disciplines ordinary language: history is not
the sentence we can tell, but the distinctions the experiment can witness.

\end{phenom}

This phenomenon is not anti-language. It is pro-accounting. The word ``after''
is useful when the record supports it. The trouble begins when grammar supplies
more order than the apparatus has earned.

\NB{Paradoxes of time travel belong here only as a caution about grammar. They
are not a Book 2 phenomenon. The physical question is not whether a clever
sentence can loop. The question is which records can be carried without
contradiction.}

\section{The Causal Doublet and Decoherence}
\label{se:the-causal-doublet-and-decoherence}

The two-slit experiment is the cleanest doublet. A coherent source of
particles --- photons, electrons, neutrons, or atoms --- illuminates a barrier
with two slits. Behind the barrier, a detector array records where each
particle arrives. With both slits open and no detector at the slits to
record which slit each particle passed through, the arrival pattern on the
detector is the famous interference pattern: alternating bright and dark
fringes whose spacing depends on the wavelength of the particle, the slit
separation, and the distance to the detector.

The pattern is built one particle at a time. Each individual particle
arrives at a single, definite location on the detector array. But the
distribution of these arrival locations, accumulated over many particles,
shows the interference pattern. This is the experimental statement of
wave-particle duality: each particle individually appears as a localized
event, while the statistical distribution requires interference between
amplitudes from both slits.

If a detector is placed at one of the slits to record which slit each
particle passes through, the interference pattern disappears. The arrival
distribution becomes the simple sum of two single-slit patterns. The
which-path record at the slit has eliminated the interference at the
detector array.

This is the chapter's central quantum phenomenon. The two-slit setup with
no which-path detection is a \emph{causal doublet}: a state that
admissibly carries both slit-traversal histories. The two-slit setup with
which-path detection is a witnessed alternative: a state in which one
specific slit-traversal history has been recorded, and the other has
not.

The transition between these two regimes is decoherence. Formally,
consider the quantum state of a particle just after the slits. With no
which-path detection, the state is a coherent superposition:
\[
|\psi\rangle = \frac{1}{\sqrt{2}}(|\text{slit 1}\rangle + |\text{slit 2}\rangle).
\]
The two-slit pattern on the detector is the square of the sum of two
complex amplitudes, with the cross-term producing the interference
fringes.

When a which-path detector interacts with the particle at the slits, the
combined system (particle + detector) becomes entangled:
\[
|\Psi\rangle = \frac{1}{\sqrt{2}}(|\text{slit 1}\rangle |\text{detector
records 1}\rangle + |\text{slit 2}\rangle |\text{detector records 2}\rangle).
\]
The particle's reduced state, traced over the detector's degrees of
freedom, becomes a classical mixture:
\[
\rho = \frac{1}{2}(|\text{slit 1}\rangle\langle\text{slit 1}| + |\text{slit 2}\rangle\langle\text{slit 2}|).
\]
The cross-terms have vanished. The detector's record of which slit was
used has destroyed the interference. The arrival distribution at the
detector array is the classical sum of two single-slit patterns.

This is decoherence. The transition from the pure coherent state to the
classical mixture is not a discontinuous collapse; it is the algebraic
consequence of tracing over an entangled environment. The environment
(the which-path detector) carries the information that distinguishes
the two slit-traversal histories. Once the environment carries that
information, interference is no longer available to the reduced
particle state.

In real experiments, the environment is rarely a single dedicated
which-path detector. More commonly, it is the laboratory's thermal
photons, scattered gas molecules, or stray electromagnetic fields ---
any degree of freedom that can interact with the particle in a way
that depends on which slit was used. A dust grain in the path is enough.
A few photons of thermal radiation are enough. The interference
pattern is fragile: it disappears as soon as any environmental degree
of freedom carries a distinguishable record of the slit.

The decoherence rate $\gamma$ quantifies this. A common model gives
$\gamma = \lambda^2 / 2D$, where $\lambda$ is the de Broglie wavelength
of the particle and $D$ is a diffusion-like constant characterizing
the environmental coupling. The interference pattern's coherence
decays as $e^{-\gamma t}$ over time $t$. For a small molecule ($\lambda
\sim 10^{-10}\,\mathrm{m}$) in a typical laboratory environment, the
decoherence time is microseconds. For a macroscopic object ($\lambda
\sim 10^{-20}\,\mathrm{m}$), the decoherence time is far shorter than
any conceivable experimental observation. This is why macroscopic
objects do not exhibit quantum interference: by the time any
measurement could be made, the environment has decohered the
superposition.

The exact form of $\gamma$ depends on the specific environmental model
(scattering cross-section, blackbody temperature, coupling Hamiltonian),
but the structural lesson is stable across models: the environment
becomes a witness. The which-path information leaves the coherent
particle state and enters the surroundings; once that transfer has
occurred, interference is no longer available to the reduced system.

\begin{phenom}{Causal Doublet}
\label{ph:causal-doublet}

\PhStatement A quantum system can carry multiple admissible histories as one
state until an interaction records which history is available to later
measurement.

\PhOrigin Two-slit and interferometer experiments exhibit interference between
alternatives that would be mutually exclusive as classical records.

\PhObservation With no which-path record, amplitudes from both slits contribute
to the observed pattern. With a which-path record, the interference terms are
suppressed.

\PhConstraint The apparatus may not compare alternatives as classical facts
until a witness has carried the distinction out of the coherent system.

\PhConsequence Decoherence is record formation. It is the point at which the
doublet stops being an admissible joint history and becomes an ordinary
ledger entry.

\end{phenom}

The word ``collapse'' can be useful, but it should not be allowed to hide
the measurement structure. Collapse is what the record looks like after
witnessing. Before that, the alternatives are transported together.
After that, the environment has made their distinction available to
later comparison.

This is why \lean{WITNESSED} is the right code anchor for the chapter. A
witness is not an opinion. It is a process that lets a later ledger entry
say: this distinction was carried.

The framework extends. Any quantum system in a superposition of two
distinguishable states is a causal doublet. The doublet can be carried
indefinitely if and only if no environmental degree of freedom acquires
which-state information. In practice, this requires extreme isolation:
ultra-high vacuum to remove scattering, cryogenic temperatures to
suppress thermal photons, electromagnetic shielding to suppress stray
fields. Quantum computing's central engineering challenge is to maintain
such isolation for long enough to perform multiple operations on the
qubits before the environment witnesses the state.

Decoherence is not a mistake or a defect. It is the structural mechanism
by which classical behavior emerges from quantum substrates. Without
decoherence, every superposition would persist indefinitely and every
measurement would face the paradox of recording two contradictory
outcomes simultaneously. With decoherence, the environment continuously
selects the basis in which records are admissible (the pointer basis),
and quantum systems coupled to that environment converge to behaviors
that match classical expectations at the macroscopic scale.

\section{Spooky Action at a Distance}
\label{se:spooky-action-at-a-distance}

Entanglement makes the doublet nonlocal. Two photons can be prepared in a joint
polarization state such that neither photon alone has a complete classical
description, but the pair has a strong correlation. Measure one and the joint
record constrains the other.

Bell's theorem gives this bite. No local hidden-variable model can reproduce
all the correlations predicted by quantum mechanics. Aspect-type experiments
then test the correlations with separated polarizer settings. The observed
violations of Bell inequalities show that the joint constraint is not a local
classical ledger waiting to be opened.

\begin{phenom}{Bell-Aspect Effect}
\label{ph:bell-aspect-effect}

\PhStatement Entangled systems can preserve a joint admissibility constraint
whose measured correlations cannot be reproduced by local hidden variables.

\PhOrigin Bell derived inequalities obeyed by local hidden-variable theories;
Aspect and later experiments observed violations using entangled photons.

\PhObservation Separated polarization measurements violate Bell inequalities
while still not permitting controllable faster-than-light signaling.

\PhConstraint The measurement may treat the pair as one carried record, but it
may not convert correlation into a usable superluminal message.

\PhConsequence Quantum transport preserves admissibility by carrying joint
constraints rather than prewritten local facts.

\end{phenom}

The old phrase ``spooky action at a distance'' is good theater and bad
mechanics if taken literally. Nothing here gives the experimenter a knob for
sending a message faster than light. The point is subtler: the admissible
record is joint before it is local. The comparison closes only when the
separated results are brought into one ledger.

This continues the theme of Chapter 6 by breaking its classical expectation.
There is still a continuation rule. There is still burden. But the record being
continued is not one path in ordinary space. It is a joint state whose
distinctions become classical only through witnessing.

\section{Davisson--Germer}
\label{se:davisson-germer}

Davisson and Germer fired electrons at a nickel crystal and observed
diffraction. The result made de Broglie's matter waves experimentally concrete:
electrons, not only light, can produce interference patterns.

For an electron with momentum $p$, the de Broglie wavelength is

\[
\lambda = \frac{h}{p}.
\label{eq:chapter7-de-broglie}
\]

At the right energy and crystal angle, the scattered electrons form peaks just
as waves scattered by a lattice should. A particle account that assigns each
electron a definite unseen path through the apparatus misses the phenomenon.

\begin{phenom}{Davisson-Germer Effect}
\label{ph:davisson-germer-effect}

\PhStatement Matter can carry wave-like admissibility, so a particle record at
the detector may be produced by interference among alternatives not separately
witnessed.

\PhOrigin The Davisson-Germer experiment observed electron diffraction from a
nickel crystal, confirming de Broglie's wavelength relation for matter.

\PhObservation Electrons scattered from the crystal show angular intensity
peaks consistent with wave diffraction rather than with independent classical
projectiles.

\PhConstraint The apparatus may not assign a definite intermediate path to an
electron when the observed record depends on coherent alternatives.

\PhConsequence Quantum measurement treats arrival events as records produced by
admissible path structure, not by hidden classical itinerary.

\end{phenom}

This is why the path integral is so natural as language, even when the chapter
does not need to develop it formally. The detector click is a finite event.
The amplitude that produces the statistics carries many admissible histories.
The click closes one record without licensing a retroactive classical path.

In the book's vocabulary, the electron's arrival is \lean{WITNESSED}. The
intermediate alternatives remain part of the carried process, not a private
movie playing behind the apparatus.

\section{Qubit Decoherence}
\label{se:qubit-decoherence}

A superconducting qubit makes decoherence painfully practical. The transmon
qubit, the dominant superconducting qubit architecture in 2025, consists of
a Josephson junction shunted by a large capacitor, fabricated on a silicon
or sapphire substrate, cooled to approximately $15\,\mathrm{mK}$ in a
dilution refrigerator. At this temperature, the thermal energy
$k_B T \approx 1{.}3 \times 10^{-25}\,\mathrm{J}$ is small enough that the
two lowest energy levels of the qubit can be addressed selectively by
microwave pulses without thermal excitation of higher levels.

The qubit's two-level structure is treated as a logical qubit. The ground
state $|0\rangle$ and the first excited state $|1\rangle$ form the
computational basis. A superposition $\alpha|0\rangle + \beta|1\rangle$ can
be prepared by a single microwave pulse of calibrated duration and phase.
Once prepared, the superposition is a causal doublet: the qubit
admissibly carries both basis states until the environment records which
one.

The environment for a superconducting qubit is everything outside the
Josephson junction and its immediate surroundings: the control wires
that bring microwave signals from room temperature down to the chip, the
substrate it sits on, the dielectric of the capacitor, the surrounding
electromagnetic field, and any defects in the materials. Each of these
provides a coupling channel through which environmental degrees of
freedom can acquire which-state information about the qubit.

Engineers describe the decoherence with two characteristic times: $T_1$
and $T_2$.

$T_1$ is the energy relaxation time. It quantifies how long the qubit
stays in $|1\rangle$ before decaying to $|0\rangle$ by emitting energy
into the environment. For a transmon at $15\,\mathrm{mK}$, $T_1$ is
typically $50$ to $200\,\mu\mathrm{s}$. The decay rate $1/T_1$ is
determined by the spectral density of the environmental noise at the
qubit's transition frequency (typically $4$ to $7\,\mathrm{GHz}$):
\[
\frac{1}{T_1} \propto S(\omega_{01}),
\]
where $S(\omega)$ is the environmental noise spectrum at frequency
$\omega$ and $\omega_{01}$ is the qubit's transition frequency in
radians per second.

$T_2$ is the dephasing time. It quantifies how long the relative phase
between $|0\rangle$ and $|1\rangle$ components of a superposition
remains coherent. Dephasing has two sources: relaxation events (which
also disturb phase), contributing a rate of $1/(2T_1)$; and pure
dephasing $T_\phi$, contributing a rate of $1/T_\phi$. The total
dephasing rate is
\[
\frac{1}{T_2} = \frac{1}{2T_1} + \frac{1}{T_\phi}.
\]
For a high-quality transmon, $T_2$ can approach $2T_1$ when pure
dephasing is suppressed. Typical values are $T_2$ in the range $50$ to
$300\,\mu\mathrm{s}$.

These are extraordinary numbers for room-temperature engineering
intuition: the qubit stays coherent for hundreds of microseconds, during
which it can be operated on by thousands of microwave pulses. But by
the standards of macroscopic objects, the qubit is fragile. A grain of
salt sitting on a laboratory bench has coherence time effectively zero
on the same scales: every interaction with the surrounding air, every
photon of thermal radiation, every vibration of the molecules, decoheres
the salt grain immediately. The qubit's hundred microseconds of coherence
is achieved by extreme isolation: cryogenic temperatures, ultra-high
vacuum, electromagnetic shielding, and careful materials engineering.

The structural framework treats $T_2$ as the chapter's quantitative
indicator of when the environment becomes a witness. For elapsed time
$t \ll T_2$, the qubit is a coherent doublet; for $t \gg T_2$, the
environment has acquired enough which-state information to render the
qubit classical. Quantum algorithms must complete within $T_2$ to use
the qubit's quantum properties; longer algorithms require quantum
error correction, which uses additional physical qubits to encode each
logical qubit and continuously monitors them for environmental
disturbances.

\begin{phenom}{Qubit Decoherence}
\label{ph:qubit-decoherence}

\PhStatement A quantum bit remains useful only while the environment fails to
carry enough information to distinguish the alternatives in its superposition.

\PhOrigin Superconducting, trapped-ion, and other qubit platforms measure
coherence times by preparing states, waiting, and testing how much phase
information remains.

\PhObservation Interactions with the environment reduce interference visibility
and eventually make the qubit behave like an ordinary classical mixture.

\PhConstraint The apparatus may compute with a doublet only while shielding,
control, and error correction prevent premature witnessing.

\PhConsequence Quantum information is not disembodied abstraction. It is a
physical record whose admissibility must be protected from unwanted witnesses.

\end{phenom}

The chapter can now answer its question. Information moves between histories
by remaining unclosed where alternatives must interfere and by closing only
where a witness is admitted. A quantum computer is therefore not magic
parallelism. It is disciplined delay: keep the record from becoming classical
until the right comparison has been performed.

The discipline is engineering-specific. The transmon qubit's $T_2$ improves
roughly an order of magnitude per decade: from $1\,\mu\mathrm{s}$ in 2005
to $100\,\mu\mathrm{s}$ in 2015 to over $1\,\mathrm{ms}$ in 2025 for the
best devices. Each order-of-magnitude improvement has required addressing
specific decoherence sources: better dielectrics in the capacitor, cleaner
fabrication processes, better isolation from control lines, lower-noise
amplifiers for readout. The progression is the chapter's framework at
work: each engineering improvement removes one channel by which the
environment was acquiring which-state information.

Other qubit platforms (trapped ions, neutral atoms, nitrogen-vacancy
centers in diamond, topological qubits) face analogous engineering
problems. Each platform has its own dominant decoherence channel, and
each platform's progress measures how well that channel can be
suppressed. The structural framework applies to all: a qubit is
admissible as a doublet only while the environment fails to witness
which alternative is occupied; the engineering task is to extend the
duration of that admissibility window.

\begin{bridge}

The chapter's question was how information can move between histories while
preserving admissibility.

The answer is witnessed transport. Interferometers show that syntactic order is
not physical order. Doublets remain coherent until an environment records their
difference. Entanglement preserves joint constraints that are stronger than
local hidden ledgers. Davisson-Germer shows matter carrying path alternatives.
Qubit decoherence makes unwanted witnessing an engineering limit.

The next question is forced by relativity. If records can be carried, witnessed,
and compared, what reference structure lets different observers agree that
their counts belong to one physical record?

\end{bridge}

\begin{coda}{A Sealed Routing Room for Diplomatic Pouches}

The routing room is windowless. Pouches arrive from embassies with seals,
numbers, and destination tags. Clerks do not open them. They scan the numbers,
place the pouches in trays, transfer the trays to couriers, and log each
handoff.

The seal is the rule. As long as it remains unbroken, the room can carry the
pouch without knowing its contents. The pouch has a route, but the contents do
not become part of the routing record. If a clerk opens the pouch, the whole
status changes. The chain of custody must now include the inspection. The
contents have been witnessed.

Two pouches may arrive from the same office with related seals. Individually,
each pouch says little. Together, once delivered and compared, they may reveal
a paired instruction: open this one only if that one arrived, discard this key
unless that packet was accepted, treat both as valid only when their numbers
match. The correlation is carried by the pair, not by either pouch alone.

No clerk can use this to send an illegal message through the room. The pouches
move by ordinary routes. The comparison happens when the records meet. But the
constraint was real before the comparison. It was sealed into the joint
procedure.

The routing room performs the chapter's structure. The unopened pouch is the
coherent system. The seal is admissibility. The log is the ledger. The clerk
who opens the pouch is the witness. The pair of correlated pouches is
entanglement. The lost or smudged seal is decoherence. The chain-of-custody
certificate is \lean{WITNESSED}.

The room is not a metaphor for ignorance alone. If it were merely ignorance,
opening the pouch would reveal a preexisting classical itinerary. Instead,
opening the pouch changes what the institution can say. Before inspection, the
contents were deliberately excluded from the routing record. After inspection,
the contents are part of the record and every later action must account for
that fact.

Quantum mechanics has the same administrative severity. The experiment may
carry alternatives without recording which one occurred. It may also create a
record that ends their interference. But it may not have both at once.

At closing time, the routing supervisor reconciles the ledger. Some pouches
were delivered unopened. Some were inspected. Some were delayed because their
partner pouch had not arrived. The day is not summarized by a single path for
each pouch independent of the rules. It is summarized by the admissible records
the room actually carried.

The seal was not decoration. It was physics in paper form.

\end{coda}
